\documentclass[11pt]{article}
\usepackage[paperwidth=1800bp, paperheight=1125bp, margin=40bp]{geometry}
\usepackage[utf8]{inputenc}
\usepackage{amsmath}
\usepackage{microtype}
\usepackage{multicol}
\usepackage{booktabs}
\usepackage{xcolor}
\usepackage{enumitem}
\usepackage{tabularx}
\usepackage{hyperref}

% Definición de Paleta Formal Tema Oscuro
\definecolor{bgdark}{HTML}{0B0F17}       % Obsidian Navy
\definecolor{primary}{HTML}{38BDF8}      % Sky Blue 400
\definecolor{primarylight}{HTML}{7DD3FC} % Sky Blue 300
\definecolor{textprimary}{HTML}{F8FAFC}  % Slate 50
\definecolor{textsecondary}{HTML}{CBD5E1}% Slate 300
\definecolor{textmuted}{HTML}{94A3B8}    % Slate 400
\definecolor{border}{HTML}{27354F}       % Slate Border
\definecolor{boxbg}{HTML}{131A29}        % Midnight Slate Panel

\setlength{\parindent}{0pt}
\setlength{\parskip}{5pt}
\setlength{\columnsep}{30bp}

\hypersetup{
    colorlinks=true,
    linkcolor=primary,
    urlcolor=primarylight
}

\begin{document}
\pagecolor{bgdark}
\color{textprimary}
\pagestyle{empty}

% ENCABEZADO INSTITUCIONAL FORMAL (DARK)
\begin{center}
    {\large\scshape\bfseries\color{primary} Proyecto de Grado en Ingeniería de Sistemas Informáticos \quad $\bullet$ \quad Universidad Católica Boliviana / Tigo SOC}\\[4bp]
    {\LARGE\bfseries\color{textprimary} MANUAL TÉCNICO Y COMPILACIÓN DE MOCKUPS MULTIPLATAFORMA}\\[4bp]
    {\normalsize\color{textsecondary} Plataforma de Detección y Mitigación de Ataques de Red en Tiempo Real \quad $\bullet$ \quad Red ISP Carrier AS27839}\\[2bp]
    {\small\color{textmuted} Postulante: Jorge Esteban Medrano \quad $\bullet$ \quad Fecha: Octubre 2026 \quad $\bullet$ \quad Entregable Maestro: Láminas Todo-en-Uno (Web, Tablet y Móvil)}
\end{center}

\vspace{-4bp}
\noindent\textcolor{primary}{\rule{\textwidth}{1.5bp}}
\vspace{4bp}

\begin{multicols}{3}

% COLUMNA 1: INSTRUCCIONES, OBJETIVO Y MARCO
\section*{\color{primary}\Large 1. Instrucciones y Objetivo}

\noindent\colorbox{boxbg}{%
\parbox{0.96\linewidth}{%
\textbf{\color{primary}Instrucciones:}\\[2bp]
\small\color{textsecondary} Como parte del proyecto de grado en Ingeniería de Sistemas Informáticos, se solicita diseñar los mockups que incluya tanto la versión web como la versión móvil (De acuerdo a cada proyecto).
}%
}

\vspace{6bp}

\noindent\colorbox{boxbg}{%
\parbox{0.96\linewidth}{%
\textbf{\color{primary}El objetivo:}\\[2bp]
\small\color{textsecondary} El objetivo es que los mockups reflejen claramente la funcionalidad principal del sistema, la experiencia de usuario, y la adaptación a diferentes dispositivos. Se debe mostrar cómo interactúan los usuarios con el sistema y cómo se garantizan los objetivos funcionales del mismo en ambas plataformas.
}%
}

\vspace{6bp}
\subsection*{\color{primarylight}Indicaciones Metodológicas}
\begin{itemize}[leftmargin=12bp, itemsep=2bp]
    \item \textbf{\color{textprimary}Diseña mockups para la interfaz web y móvil:} Láminas que unifican Web Desktop (1440$\times$900), Tablet (1024$\times$768) y Móvil PWA (390$\times$844) en una misma página de 2400$\times$1500 px.
    \item \textbf{\color{textprimary}Pantallas y flujos clave:} 10 pantallas críticas del SOC modeladas con interacciones reales y telemetría de kernel Linux.
    \item \textbf{\color{textprimary}Coherencia funcional:} La lógica de negocio, contratos de API REST y respuestas eBPF/XDP son idénticos en todas las plataformas.
    \item \textbf{\color{textprimary}Herramientas y técnicas adecuadas:} Modelado gráfico vectorial SVG, compilación automatizada y tipografía tabular sin distorsión.
    \item \textbf{\color{textprimary}Decisiones de diseño:} Enfoque riguroso en usabilidad, ergonomía y prevención de errores en salas de operaciones 24/7.
\end{itemize}

\columnbreak

% COLUMNA 2: DECISIONES DE DISEÑO, USABILIDAD Y FUNCIONALIDAD
\section*{\color{primary}\Large 2. Usabilidad y Funcionalidad}

\subsection*{\color{primarylight}Erradicación de Anti-Patrones UI/UX de IA}
\begin{itemize}[leftmargin=12bp, itemsep=2bp]
    \item \textbf{\color{textprimary}Cero Emojis:} Eliminación absoluta de emojis informales en estados, tablas y botones. Se emplea iconografía técnica vectorial SVG pura y texto explícito.
    \item \textbf{\color{textprimary}Textos y Especificaciones Semi-Planos:} Contenedores con bordes uniformes de 1\,px (\emph{sin barras gruesas a la izquierda}). Eliminación del anidamiento excesivo (\emph{cardocalypse}).
    \item \textbf{\color{textprimary}Radios Concéntricos Calibrados:} Geometría armónica según $\text{radio}_{\text{interior}} = \max(0, \text{radio}_{\text{exterior}} - \text{padding})$.
    \item \textbf{\color{textprimary}Tipografía Tabular de Red:} Fuentes sans-serif limpias y tipografías monoespaciadas con cifras tabulares para IP, CIDR, hashes y pps.
\end{itemize}

\subsection*{\color{primarylight}Pilares de Seguridad Reforzada}
\begin{itemize}[leftmargin=12bp, itemsep=2bp]
    \item \textbf{\color{textprimary}Sesión Inmutable en Header:} Eliminación de cambio arbitrario de cuenta. Despliegue fijo de \texttt{Carlos Mendoza [admin\_soc]}, rol \texttt{SOC\_ADMIN}, estado TLS 1.3 y botón de \texttt{[Cerrar Sesión]}.
    \item \textbf{\color{textprimary}Rotación Obligatoria de Contraseña:} Detección de clave temporal por defecto (\texttt{admin}). Bloqueo de acceso hasta definir clave con medidor de entropía (94 bits) y checklist ISO 27001.
    \item \textbf{\color{textprimary}Autenticación Corporativa de 3 Factores:} Correo SSO institucional, casillas numéricas 2FA/TOTP sincronizadas con app móvil y rate limiting (3 intentos).
    \item \textbf{\color{textprimary}Sandbox ACL en Subred Aislada:} Subred virtualizada \texttt{10.200.0.0/24 (VLAN 999)} con inyector sintético, consola eBPF/iptables y promoción auditada.
\end{itemize}

\columnbreak

% COLUMNA 3: ÍNDICE DE LAS 10 PÁGINAS Y MATRIZ DE ENTREGABLES
\section*{\color{primary}\Large 3. Índice de Pantallas del SOC}

\color{textsecondary}
\begin{tabularx}{\linewidth}{lX}
\toprule
\textbf{\color{primary}Página} & \textbf{\color{primary}Funcionalidad y Controles Clave} \\
\midrule
\textbf{\color{textprimary}Página 1} & \textbf{\color{textprimary}Login de Alta Seguridad}\\[-1bp]
\texttt{01\_login} & SSO Email, 2FA TOTP 6 dígitos, mTLS 1.3 y protección contra fuerza bruta (Redis rate-limit). \\
\midrule
\textbf{\color{textprimary}Página 2} & \textbf{\color{textprimary}Rotación Obligatoria de Contraseña}\\[-1bp]
\texttt{01b\_cambio} & Medidor de entropía dinámica (94 bits) y validación de 5 reglas criptográficas ISO 27001. \\
\midrule
\textbf{\color{textprimary}Página 3} & \textbf{\color{textprimary}Dashboard Central del SOC}\\[-1bp]
\texttt{02\_dashboard} & 4 KPIs de transporte (42.8 Gbps), curva de ancho de banda y mitigación en kernel. \\
\midrule
\textbf{\color{textprimary}Página 4} & \textbf{\color{textprimary}Topología de Red Core GNS3}\\[-1bp]
\texttt{03\_topology} & Mapa AS27839, enlaces 10G/40G e inspector lateral de interfaces eth0-eth3. \\
\midrule
\textbf{\color{textprimary}Página 5} & \textbf{\color{textprimary}Métricas ML (LightGBM)}\\[-1bp]
\texttt{04\_monitoring} & Inferencia 0.42 ms, precisión 99.42\%, ranking de variables y matriz de confusión. \\
\midrule
\textbf{\color{textprimary}Página 6} & \textbf{\color{textprimary}Inspección eBPF en Tiempo Real}\\[-1bp]
\texttt{05\_realtime} & Captura de 14.2k pkts/s, filtros BPF y deserializador hexadecimal/ASCII. \\
\midrule
\textbf{\color{textprimary}Página 7} & \textbf{\color{textprimary}Auditoría Forense y Alertas}\\[-1bp]
\texttt{06\_alerts} & Registro de incidentes, severidades, firma SHA-256 y exportación pericial. \\
\midrule
\textbf{\color{textprimary}Página 8} & \textbf{\color{textprimary}Control de Blast Radius}\\[-1bp]
\texttt{07\_blast\_rad} & Cuarentena preventiva, temporizador de rollback 300s y re-autenticación FIDO2. \\
\midrule
\textbf{\color{textprimary}Página 9} & \textbf{\color{textprimary}Gobernanza RBAC Zero-Trust}\\[-1bp]
\texttt{08\_user\_mgmt} & Administración de operadores, enforcement 2FA al 100\% y matriz de privilegios. \\
\midrule
\textbf{\color{textprimary}Página 10} & \textbf{\color{textprimary}Sandbox ACL en Subred Aislada}\\[-1bp]
\texttt{09\_sandbox} & Simulación en VLAN 999 (\texttt{10.200.0.0/24}), métricas de mitigación y promoción. \\
\bottomrule
\end{tabularx}

\end{multicols}

\vspace{6bp}
\noindent\textcolor{border}{\rule{\textwidth}{0.8bp}}
\vspace{2bp}
\begin{center}
    {\footnotesize\color{textmuted}\textsc{Compilación Maestra Multi-Dispositivo} \quad $\bullet$ \quad Todas las láminas presentan Web Desktop, Tablet y Móvil integradas con el mismo tamaño uniforme \quad $\bullet$ \quad Cero Emojis $\bullet$ Estilo Semi-Plano}
\end{center}

\end{document}
